\documentclass[11pt,a4paper]{article}
\usepackage[utf8]{inputenc}
\usepackage{amsmath,amssymb,amsthm}
\usepackage{listings}
\usepackage{xcolor}
\usepackage[margin=1in]{geometry}
\usepackage{hyperref}

\newtheorem{theorem}{Theorem}
\newtheorem{lemma}[theorem]{Lemma}

\lstset{
  language=C++,
  basicstyle=\ttfamily\small,
  keywordstyle=\color{blue}\bfseries,
  commentstyle=\color{green!60!black},
  stringstyle=\color{red},
  numbers=left,
  numberstyle=\tiny\color{gray},
  breaklines=true,
  frame=single,
  tabsize=4,
  showstringspaces=false
}

\title{IOI 2004 -- Hermes}
\author{Solution}
\date{}

\begin{document}
\maketitle

\section{Problem Statement Summary}
Hermes starts at the origin and must visit $N$ events in order. Event
$i$ is at position $p_i$ on axis $d_i \in \{X, Y\}$. He moves along
both axes \emph{simultaneously}: the time to go from $(x_1, y_1)$ to
$(x_2, y_2)$ is the Chebyshev distance
$\max(|x_2 - x_1|, |y_2 - y_1|)$. Minimize total travel time.

\section{Solution: Interval DP on the Free Axis}

\subsection{Key Observation}
After visiting event $i$ (which lies on axis $d_i$), the coordinate on
axis $d_i$ is fixed at $p_i$. The coordinate on the \emph{other} axis
is free---Hermes can choose it, subject to the time budget. The optimal
strategy tracks an \emph{interval} of reachable positions on the free
axis.

\subsection{State}
After processing event $i$, maintain:
\begin{itemize}
  \item $\mathrm{cost}$: total time spent so far (minimum possible).
  \item $[\mathrm{lo}, \mathrm{hi}]$: the interval of positions on the
        free axis that are achievable at cost $\mathrm{cost}$.
\end{itemize}

\subsection{Transitions}
\paragraph{Same-axis transition} ($d_i = d_{i-1}$).
The event axis moves from $p_{i-1}$ to $p_i$; this takes at least
$\Delta = |p_i - p_{i-1}|$ time. During that time, the free axis can
drift by up to $\Delta$ in either direction:
\[
  \mathrm{cost} \mathrel{+}= \Delta, \qquad
  \mathrm{lo} \mathrel{-}= \Delta, \qquad
  \mathrm{hi} \mathrel{+}= \Delta.
\]

\paragraph{Cross-axis transition} ($d_i \ne d_{i-1}$).
The old free axis (now axis $d_i$) had position in $[\mathrm{lo}, \mathrm{hi}]$.
We must move it to $p_i$. The minimum extra time is:
\[
  \delta = \max\bigl(0,\; p_i - \mathrm{hi},\; \mathrm{lo} - p_i\bigr)
  = \operatorname{dist}(p_i, [\mathrm{lo}, \mathrm{hi}]).
\]
The new free axis (the old event axis $d_{i-1}$, fixed at $p_{i-1}$)
can drift by up to $\delta$ during this extra time:
\[
  \mathrm{cost} \mathrel{+}= \delta, \qquad
  \mathrm{lo} = p_{i-1} - \delta, \qquad
  \mathrm{hi} = p_{i-1} + \delta.
\]

\begin{lemma}
The interval representation is exact: every position in
$[\mathrm{lo}, \mathrm{hi}]$ is achievable at the stated cost, and no
position outside is achievable at lower cost.
\end{lemma}
\begin{proof}
The Chebyshev distance is the maximum of independent movements on the
two axes. Since only one axis is constrained per event, the other axis
can freely absorb up to $\mathrm{cost}$ of movement. The interval
tracks the exact range of ``slack'' on the free axis.
\end{proof}

\section{C++ Implementation}

\begin{lstlisting}
#include <bits/stdc++.h>
using namespace std;

int main() {
    ios::sync_with_stdio(false);
    cin.tie(nullptr);

    int N;
    cin >> N;

    vector<int> d(N), p(N);
    for (int i = 0; i < N; i++) {
        cin >> d[i] >> p[i];
        d[i]--; // 0 = X-axis, 1 = Y-axis
    }

    // Initial move: from (0,0) to event 0 at p[0] on axis d[0].
    // Time = max(|p[0]|, |free|) >= |p[0]|; minimised at |p[0]|.
    // Free axis can be anywhere in [-|p[0]|, |p[0]|].
    long long cost = abs(p[0]);
    long long lo = -abs(p[0]);
    long long hi = abs(p[0]);

    for (int i = 1; i < N; i++) {
        if (d[i] == d[i - 1]) {
            // Same axis: mandatory cost = |p[i] - p[i-1]|
            long long dt = abs(p[i] - p[i - 1]);
            cost += dt;
            lo -= dt;
            hi += dt;
        } else {
            // Cross axis: must reach p[i] from interval [lo, hi]
            long long delta;
            if (p[i] >= lo && p[i] <= hi)
                delta = 0;
            else if (p[i] < lo)
                delta = lo - p[i];
            else
                delta = p[i] - hi;

            cost += delta;
            lo = (long long)p[i - 1] - delta;
            hi = (long long)p[i - 1] + delta;
        }
    }

    cout << cost << "\n";
    return 0;
}
\end{lstlisting}

\section{Complexity Analysis}
\begin{itemize}
  \item \textbf{Time:} $O(N)$---a single pass over the events.
  \item \textbf{Space:} $O(N)$ for input storage (the DP uses $O(1)$
        extra space).
\end{itemize}

\end{document}
